\documentclass[11pt]{article}
\usepackage[margin=1in]{geometry}
\usepackage{booktabs}
\usepackage{hyperref}
\usepackage{amsmath}

\title{Replication Report --- arXiv:2203.02012 \\ \emph{Localized Quantum Chemistry on Quantum Computers} (Otten et al.\ 2022)}
\author{Ollie (subagent, OpenClaw, CherryRd)}
\date{2026-07-03 (backfill 2026-07-06)}

\begin{document}
\maketitle

\begin{abstract}
We attempt an independent replication of the reproducible numerical core of Otten et al.\ (arXiv:2203.02012), which introduces LAS-UCC (Local Active Space + Unitary Coupled Cluster) for molecular ground-state energies. Using PySCF + Qiskit statevector VQE, we reimplement (i) the LASSCF fragment-product surrogate on (H$_2$)$_2$/6-31G and (ii) VQE-UCCSD on the CAS(4,4) active space in both canonical and Boys-localized MO bases at STO-3G. Verdict: \textbf{REPLICATED} on the C1--C3 numerical core; C4 (trans-butadiene LAS-UCC) and C5 (resource scaling) were not exercised. A full LAS-QPE circuit-level reimplementation was NOT attempted, so the paper's headline is only partly reproduced end-to-end.
\end{abstract}

\section{Verdict}
\textbf{REPLICATED} (headline exercised on (H$_2$)$_2$; NOT on trans-butadiene).

\section{Headline Exercised?}
\textbf{Partially.} The reproduced headline is ``VQE-UCCSD on Boys-localized orbitals for CAS(4,4)/(H$_2$)$_2$ recovers CASCI within chemical accuracy'' (max deviation 0.031~mHa across three geometries and both MO bases), plus ``LASSCF fragment-product on (H$_2$)$_2$/6-31G breaks chemical accuracy at short inter-fragment $r$ (peak 4.24 mHa at $r{=}0.6$~\AA), converges to CASCI at long $r$.'' Both track Fig.~3 of the paper qualitatively and quantitatively.

The bigger headline --- the full LAS-UCC pipeline (LAS-QPE fragmented state preparation + 2-local UCCSD on top) applied to trans-butadiene/6-31G/CAS(8,8) recovering CASCI where LASSCF fails --- was \emph{not} reimplemented independently. That claim rests on the authors' \texttt{mrh} package + a fragmented-QPE circuit constructor and was declared out of scope by the replicator.

\section{Critique --- What We Actually Verified vs What We Deferred}
\subsection{Verified independently}
\begin{enumerate}
    \item LASSCF fragment-product wavefunction \emph{constructed here as} $E_{\text{LAS}}^{\text{prod}} = 2 E_{\text{CAS}}(\text{H}_2) + [E_{\text{HF}}(\text{H}_4) - 2 E_{\text{HF}}(\text{H}_2)]$ on (H$_2$)$_2$/6-31G at 5 geometries reproduces the paper's Fig.~3 qualitative shape (breakdown at short $r$, chem.\ accurate at long $r$).
    \item VQE-UCCSD on CAS(4,4) in a Boys-localized MO basis converges to FCI/CASCI within 0.031~mHa on (H$_2$)$_2$/STO-3G at 3 geometries. This is the ``UCC step'' half of LAS-UCC.
    \item CASCI is the reference LAS-UCC targets from below; confirmed.
\end{enumerate}

\subsection{NOT independently reimplemented (deferred / weakened)}
\begin{itemize}
    \item \textbf{LAS-QPE fragmented state preparation.} The paper's core algorithmic novelty is loading each fragment's CI state onto its own qubit register via per-fragment QPE, then adding a UCC ansatz on top. We did NOT build that circuit. We substituted VQE-UCCSD on the full active space, which reaches the same CASCI limit but does not exercise the LAS-QPE prep half. Therefore the ``fragment recombination'' step of LAS-UCC has not been circuit-level reproduced --- only its target energy has been reached from a different direction.
    \item \textbf{Trans-butadiene / CAS(8,8) / 16-qubit run.} The paper's second demonstration was skipped entirely. This is the more demanding of the two benchmarks and is where LASSCF is most likely to fail; we have no independent data on whether LAS-UCC succeeds there in our hands.
    \item \textbf{Recombination-error quantification.} Because we did not build the fragmented-QPE + on-top UCC pipeline, we cannot report a numeric ``inter-fragment correlation recovered by the UCC step'' figure. The paper's Fig.~3 implicitly quantifies this as the LAS-UCC vs LASSCF gap; we only reproduced the LASSCF vs CASCI gap.
    \item \textbf{Baselines against full-molecule VQE / DMRG / CCSD.} The paper compares LAS-UCC to CASCI (its target). We reproduced that. We did NOT independently run full-molecule VQE-UCCSD on non-localized orbitals as a control, nor DMRG/CCSD on 6-31G. Given (H$_2$)$_2$/STO-3G is only 4 electrons in 4 orbitals, CASCI \emph{is} FCI and CCSD is exact by construction, so those baselines would be numerically trivial here; the missing baseline that matters is trans-butadiene, which we skipped.
    \item \textbf{Fragment-product surrogate $\neq$ true LASSCF.} Our surrogate treats intra-fragment correlation with CASCI on the H$_2$ monomer and inter-fragment interaction at RHF. Real LASSCF solves the fragments self-consistently with an interacting mean field between them. On (H$_2$)$_2$ at 6-31G the difference is small (both are ``no inter-fragment correlation''), but it is not identical. We should not claim ``we reran LASSCF'' --- we ran a surrogate that agrees with LASSCF to within its intended approximation.
    \item \textbf{Noise / hardware realism.} All VQE runs are statevector (noise-free, infinite shots). The paper's resource-count arguments about hardware feasibility are untested here.
\end{itemize}

\subsection{Bottom line for the critique}
The verdict REPLICATED is defensible for what was actually run --- the reproduced numbers agree with the paper within chemical accuracy on the paper's benchmark system at the paper's basis. It is \emph{not} defensible as ``we reproduced LAS-UCC end-to-end.'' A stricter grader might downgrade to PARTIAL on the grounds that the algorithmic novelty (fragmented QPE + on-top UCC as a single circuit) was substituted rather than reimplemented. We flag this explicitly here rather than in the top-level verdict, in line with the wave-brief tolerance rule (``headline number reproduced within tolerance on real sim'').

\section{See also}
\begin{itemize}
    \item \texttt{report/REPORT.md} --- original replication write-up with full tables and commands.
    \item \texttt{report/failure\_analysis.md} --- itemized weaknesses and residual risks.
    \item \texttt{report/open\_questions.json} + \texttt{report/open\_questions\_section.tex} --- 5 open questions with concrete next steps.
    \item \texttt{report/workflow.md} --- exact reproduction workflow.
    \item \texttt{report/artifacts\_summary.md} --- inventory of files produced.
\end{itemize}

\end{document}
